\section{Derivation of the update rule}
\label{app:derivation}

Everything the source draft spends its \S II on belongs here. The main text needs
Eqs.~\ref{eq:fields}--\ref{eq:update} and the two properties of
\S\ref{sec:rule}; a reader who wants to check the derivation wants it complete and
in one place.

\paragraph{Order to present it in.}
\begin{enumerate}
  \item The likelihood. Two corruptions, marginalised: the channel flip with
    probability $\varepsilon(z) = \Phi(z)$, $z \sim \mathcal{N}(\mu_{e|r}, V_{e|r})$,
    and the encoding mismatch with $\mathbb{E}[x_i y_i] = v_p$, $v_p = \|\wv\|^{-2}$.
    Result: $L(\sigma_e \mid z, \xv, \wv) = \Phi(z)\Phi(-\wv\cdot\xv\,\sigma_e)
    + (1 - \Phi(z))\Phi(\wv\cdot\xv\,\sigma_e)$.
  \item Why $\varepsilon = \Phi(z)$ rather than a Beta posterior. The Beta choice
    is the natural one and breaks conjugacy; $\Phi$ keeps the state in the Gaussian
    family of Eq.~\ref{eq:state}. State this as the modelling decision it is.
  \item The evidence. Gaussian integrals give Eq.~\ref{eq:evidence}. Note the
    exclusive-or form explicitly --- it is the cleanest statement of what surprise
    means in this model and it is worth a sentence rather than being left to the
    reader.
  \item The maximum-entropy projection back onto Eq.~\ref{eq:state}, matching first
    and second moments. This is the step that makes the rule an instance of
    entropic dynamics and the step that makes it \emph{optimal for a pair}, which
    the introduction leans on. Give it properly.
  \item The modulation functions as log-derivatives of the evidence,
    Eq.~\ref{eq:modulation}, and the resulting update Eq.~\ref{eq:update}. One
    presentation, not three. The source draft gives the same update in three
    notations (its Eqs.~9--12, 13--16, 17--20); collapse them.
  \item The symmetry $F_w(x,y) = F_\mu(y,x)$, $F_C(x,y) = F_V(y,x)$, and why the
    scaling by $\gamma_C, \gamma_V$ is what produces it.
\end{enumerate}

\paragraph{Numerical safeguards.} The exact modulation functions diverge where
$Z \to 0$, deep in the dissonant corners, and $F_V < 0$ drives $V$ monotonically to
zero, so floors on $Z$ and $V$ are needed for a long run to survive. The covariance
update $C \to C + a\,(C\xv)(C\xv)^{\!\top}$ stays positive semi-definite only while
$1 + a\,\xv\cdot C\xv \geq 0$, and $F_C$ is unbounded below, so $a$ is clipped at
that boundary. Report the floors, report how often the clip triggers, and say
plainly that only the clip could alter a trajectory. The source draft mentions none
of this; a reviewer who tries to reimplement will hit all three.

\begin{figure}[h]
  \centering
  \includegraphics[width=\textwidth]{modulation_contours}
  \caption{The four modulation functions of Eq.~\ref{eq:modulation} over
    $(\hw, \hmu)$. The diagonal $|\hw| = |\hmu|$ is the blame-attribution crossover
    of \S\ref{sec:rule}; the sign change of $F_w$ across $\hmu = 0$ is the
    anti-learning that ablation A2 removes.
    \todo{axes are currently labelled with the unscaled fields; relabel as $\hw$,
    $\hmu$ --- this is flagged in the source draft too and never fixed.}}
  \label{fig:modulation}
\end{figure}

\todo{\texttt{modulation\_surfaces} and \texttt{modulation\_slices} are also
available. Pick one of the three for the appendix; three views of the same four
functions is more than this needs.}
